% =========================================================
% 4.4 INSTRUMENTOS DE RECOLECCIÓN DE DATOS
% El \section{} de esta sección está en el chapter.tex del capítulo;
% sus referencias van en seccion4_4.bib (misma carpeta).
% =========================================================

La recolección de información se realiza a partir de tres fuentes complementarias:

\begin{itemize}
    \item \textbf{Datos oficiales.} Se utilizan los datos de incidencia de delitos violentos (\textit{especificar tipos}) de la Fiscalía General de Justicia, así como el portal de datos abiertos, correspondientes al periodo de enero de 2024 a \textit{agosto} de 2026.
    \item \textbf{Imágenes de Google Street View.} Utilizadas para la extracción de variables del entorno construido, corresponden a enero de 2025. Se asume que las características físicas del entorno se mantienen relativamente estables durante el periodo de análisis.
    \item \textbf{Datos crowdsourced.} Se plantea la utilización de fuentes alternas de información como lo pueden ser reportes en redes sociales, encuestas y datos extraídos del mismo prototipo de aplicación móvil.
\end{itemize}
